\section{Demonstration Scenarios}
\label{sec:demonstration}
\begin{figure*}[t] % 浣跨敤 figure* 鍙互璺ㄥ弻鏍忔樉绀猴紙濡傛灉鏄崟鏍忔ā鏉垮幓鎺?锛?  \centering
  \includegraphics[width=0.9\textwidth]{Fig/ui_v3.pdf} 
  \caption{A running example of \modora}
  \label{fig:gui}
\end{figure*}

This section demonstrates how \modora turns semi-structured documents into verifiable knowledge assets (Figure~\ref{fig:gui}).

\runin{Hierarchical Parsing and Inspection} After ingestion, \modora executes an asynchronous pipeline (\textit{OCR/fallback} $\rightarrow$ \textit{component aggregation} $\rightarrow$ \textit{tree construction} $\rightarrow$ \textit{metadata generation}) and stores intermediate artifacts (e.g., OCR blocks, component pack, tree JSON). Users inspect and refine this structure through the \textit{Tree Visualizer} (Figure~\ref{fig:gui}-\circled{2}).

\runin{Verified Multimodal QA} Users issue natural-language queries across documents, and \modora performs query-adaptive retrieval (location or semantic+vector). With \textit{Interactive Grounding} (Figure~\ref{fig:gui}-\circled{5}), citation clicks highlight source regions in the original PDF layout (Figure~\ref{fig:gui}-\circled{6}). Each evidence item keeps retriever provenance (location/semantic/vector), page index, and bounding boxes for transparent validation.

\runin{Human-in-the-Loop Optimization} The \textit{Node Heatmap} (Figure~\ref{fig:gui}-\circled{3}) visualizes retrieval impact, while users can refine the tree through direct node editing or natural-language restructuring instructions. Indexed CCTrees are archived in the \textit{Knowledge Base View} (Figure~\ref{fig:gui}-\circled{4}) for cross-session reuse, together with auto-generated document tags and structural statistics.
